\begin{textsample}{POS Dim 1 – ro – Score 40.00 – ro/ro\_em\_45.txt}  \label{ex:f1_pos_148}
2.4.1. \textbf{Categorias} e a Articulação entre os Componentes da Área A área de Ciências Humanas e Sociais aplicadas traz \textbf{uma} especificidade \textbf{que} são as \textbf{categorias} como linha transversal, garantindo a consonância entre todos os componentes da área, o \textbf{que} de modo geral garante a \textbf{interdisciplinaridade} e a \textbf{transdisciplinaridade}, para \textbf{que} haja \textbf{uma} educação de forma integral, a fim de \textbf{que} haja o desenvolvimento humano, foco principal da área.
Diante desse contexto, é necessário conhecer e reconhecer as peculiaridades regionais e compreender as necessidades da juventude da \textbf{contemporaneidade} é por esse motivo principalmente a necessidade de articulação dos objetos de conhecimento dentro da área, dialogando principalmente com temas da \textbf{contemporaneidade}.
Considerando o trabalho educativo com o desenvolvimento de competências e habilidades, farão com \textbf{que} os estudantes se desenvolvam \textbf{enquanto} ser, e compreendam a importância do reconhecimento do outro, como \textbf{uma} proposta para se fazer \textbf{uma} sociedade mais \textbf{justa} e mais fraterna.
Com efeito, entende -se \textbf{que}, \textbf{uma} vez \textbf{que} o \textbf{sujeito} se conhece e se percebe no mundo, \textbf{que} percebe a presença do outro, e compreende \textbf{que} cada \textbf{um} sente a necessidade de se relacionar com o outro, surgem, portanto, os primeiros princípios de cidadania e ética, \textbf{que} serão fundamentais para a boa convivência entre os seres.
Por \textbf{conseguinte}, para fundamentar ainda mais essa relação ( Eu + Outro = Nós ), A BNCC estabelece quatro \textbf{categorias} específicas para a área de CHSA, \textbf{que} serão determinantes no Ensino Médio para o desenvolvimento do indivíduo, a saber : Tempo e Espaço ; Território e \textbf{Fronteira} ; Indivíduo, Natureza, Sociedade, Cultura e Ética ; e Política e Trabalho.
Considerando a capacidade de “ abstração ” e “ simbolização ”, essas \textbf{categorias} visam \textbf{um} aprofundamento ainda maior daquelas \textbf{que} foram trabalhadas na etapa do Ensino Fundamental.
Faz -se necessário analisar o \textbf{que} a Base Nacional Comum Curricular versa sobre cada \textbf{uma} delas.
A primeira \textbf{categoria} apresentada é Tempo e Espaço.
A princípio, podemos pensar \textbf{que} essa \textbf{categoria} \textbf{diz} respeito somente aos componentes curriculares de História e Geografia, no entanto, a Filosofia e a Sociologia também abordam essa \textbf{categoria}.
Muitos filósofos e \textbf{cientistas} ao longo da história já tentaram dar suas versões em relação aos conceitos apresentados.
Toda história da vida humana se dá dentro de \textbf{um} tempo cronológico e dentro de \textbf{um} espaço territorial.
No entanto, não devemos ter a pretensão de definir esses conceitos de \textbf{uma} forma \textbf{simples} e direta.
A definição de tempo e espaço \textbf{que} a BNCC apresenta, vai muito além daquilo \textbf{que} é comum e trivial para o ser humano. \textbf{Vejamos} o \textbf{que} o documento da BNCC \textbf{diz} sobre essa \textbf{categoria} : Tempo e espaço explicam os fenômenos nas Ciências Humanas porque permitem identificar contextos [... ].
A localização no tempo e no espaço nos permite identificar circunstâncias, tornando possível compará -las, observar as semelhanças e as diferenças, assim como as permanências e as transformações.
Definir o \textbf{que} seria o tempo é \textbf{um} desafio sobre o qual se debruçaram e se debruçam grandes pensadores de diversas áreas do conhecimento.
O tempo é matéria de reflexão na Filosofia, na Física, na Matemática, na Biologia, na História, na Sociologia e em outras áreas do saber ( BNCC, 2017, p. 563-564 ).
Por sua vez, a compreensão do espaço contempla as dimensões histórica e cultural, ultrapassando suas representações \textbf{cartográficas}.
O Espaço está associado aos arranjos dos objetos de diversas naturezas, mas também às movimentações das sociedades, nas quais ocorrem eventos, disputas, conflitos, ocupações ( ordenadas e desordenadas ) ou dominações.
No espaço ( em \textbf{um} lugar ) se dá a produção, distribuição e consumo de mercadorias.
Nele são realizados fluxos de diversas naturezas ( pessoas e objetos ) e são desenvolvidas relações de trabalho, com ritmos e velocidade variados ( BNCC, 2017, p. 563-564 ).
A segunda, é a de Território e \textbf{Fronteira}.
Ambos os conceitos parecem ser fáceis, assim como a \textbf{categoria} de Tempo e Espaço, mas não são : eles são \textbf{bastante} amplos.
Quando se fala em território e \textbf{fronteira}, identifica -se todas as mudanças \textbf{que} ocorreram e ocorrem dentro das cidades e/ou regiões e \textbf{que} denotam as marcações de \textbf{um} território ou de vários, com todas as suas formas de organizações sociais ( religiões, culturas, vestimentas, festas, tradições etc. ), \textbf{que} são diferentes de outras localidades, pois cada \textbf{uma} possui características específicas.
Cada cidade, cada região, cada país possui seus territórios e suas \textbf{fronteiras}, sendo \textbf{que} algumas delas podem ser motivos de lutas e conquistas, e de derrotas.
Foi pensando na multiplicidade de \textbf{posicionamentos} em relação a essa \textbf{categoria} \textbf{que} a BNCC se refere da seguinte forma : Território é \textbf{uma} \textbf{categoria} usualmente associada a \textbf{uma} porção da superfície terrestre sob domínio de \textbf{um} grupo e suporte para nações, estados, países.
É dele \textbf{que} provêm alimento, segurança, identidade e refúgio.
Engloba as noções de lugar, região, \textbf{fronteira} e, especialmente, os limites políticos e administrativos das cidades, estados e países, sendo, portanto, esquemas abstratos de organização da realidade.
Associa -se a ele também a ideia de poder, jurisdição, administração e soberania, dimensões \textbf{que} expressam a diversidade das relações sociais e permitem juízos analíticos ( Ibidem, p. 564-565 ). \textbf{Fronteira} também é \textbf{uma} \textbf{categoria} construída \textbf{historicamente}.
Ao expressar \textbf{uma} cultura, grupos definem \textbf{fronteiras}, formas de organização social e, por vezes, áreas de confronto com outros grupos.
A conformação dos impérios coloniais, a formação dos Estados Nacionais e os processos de globalização problematizam a discussão sobre limites culturais e \textbf{fronteiras} nacionais.
Os limites, por exemplo, entre civilização e barbárie geraram, não raro, a destruição daqueles indivíduos considerados bárbaros.
Temos aí \textbf{uma} \textbf{fronteira} sangrenta.
Povos com culturas e saberes distintos em muitos casos foram separados ou reagrupados de forma a resolver ou agravar conflitos, facilitar ou dificultar deslocamentos humanos, favorecer ou impedir a integração territorial de populações com identidades semelhantes ( Ibidem, p. 564-565 ).
A terceira \textbf{categoria}, Indivíduo, Natureza, Sociedade, Cultura e Ética é \textbf{bastante} ampla, pois envolve vários conceitos \textbf{que} são determinantes para o desenvolvimento do estudante.
Essa \textbf{categoria} oferece \textbf{uma} ampla discussão a respeito dos esclarecimentos \textbf{teóricos} tendo como base as civilizações antigas.
Essas civilizações tinham \textbf{uma} forma de organização política, econômica, social, ambiental e cultural diferente \textbf{uma} das outras.
Cada sociedade se organizava de forma independente e se estruturava conforme seus costumes.
Essa \textbf{categoria} tem em sua \textbf{centralidade} o \textbf{homem} : o \textbf{homem} \textbf{que} está inserido em \textbf{um} contexto e \textbf{que} \textbf{necessita} ser estudado e compreendido a partir de sua realidade.
No \textbf{que} \textbf{diz} respeito ao conhecimento do \textbf{homem}, os gregos propuseram \textbf{uma} ampla discussão a respeito do assunto.
O filósofo grego Sócrates indagava os transeuntes com a célebre frase, “ conhece -te a ti mesmo ”, e essa pergunta sobre quem é o \textbf{homem}, sobre a necessidade de conhecer a si mesmo, perpassou por toda história da vida humana.
Vários \textbf{pesquisadores}, de diversas áreas do conhecimento, já tentaram encontrar \textbf{uma} resposta plausível para tal pergunta.
Até \textbf{hoje}, pouco se conhece sobre o \textbf{homem}, mesmo \textbf{que} já se tenha descoberto muito sobre ele.
O \textbf{homem} é \textbf{um} ser \textbf{que} conhece, \textbf{que} é dotado de razão, de inteligência.
No entanto, o \textbf{homem} está inserido em \textbf{um} mundo \textbf{que} está em constante modificação.
Essas mudanças ocorrem em vários cenários : no campo político, econômico, social, ambiental, cultural etc.
Como já foi observado, essa \textbf{categoria} \textbf{que} iremos abordar é \textbf{bastante} ampla, pois ela envolve vários conceitos, \textbf{que} são :"Indivíduo", “ Natureza ”, “ Sociedade ”, “ Cultura ” e “ Ética ”.
Poderíamos abordar cada \textbf{um} deles de forma separada ( \textbf{que} poderia gerar até \textbf{uma} melhor compreensão ), ou poderíamos abordá -la como \textbf{um} todo, pois todos esses conceitos estão relacionados a \textbf{um} único ponto : o “ \textbf{homem} ”. \textbf{Vejamos} o \textbf{que} essa \textbf{categoria} tem a \textbf{dizer} sobre cada conceito. a ) Sobre o conceito de Indivíduo : Na modernidade, a noção de indivíduo se tornou mais complexa em razão das transformações ocorridas no âmbito das relações sociais marcadas por novos códigos culturais, concepções de individualidade e formas de organização política.
Em meio às mudanças, foram criadas condições para o debate a respeito da natureza dos seres humanos, seu papel em diferentes culturas, suas instituições e sua capacidade para a autodeterminação.
A sociedade capitalista, por exemplo, ao mesmo tempo em \textbf{que} propõe a \textbf{centralidade} de \textbf{sujeitos} iguais, constrói relações econômicas \textbf{que} produzem e reproduzem desigualdades no corpo social ( Ibidem, p. 554 ).
As diferenças e semelhanças entre os indivíduos e as sociedades foram sedimentadas ao longo do tempo e em múltiplos espaços e circunstâncias.
Procurar identificar essas diferenças e semelhanças tanto em seu grupo social ( familiar, escolar, bairro, cidade, país, etnia, religião etc. ) quanto em outros povos e sociedades constitui \textbf{uma} aprendizagem a ser garantida aos estudantes na área de Ciências Humanas e Sociais Aplicadas ( Ibidem, p. 554 ). b ) Sobre o conceito de Sociedade e Natureza : A sociedade, da qual faz parte o indivíduo, consiste em \textbf{um} grupo humano, ocupante de \textbf{um} território, com \textbf{uma} forma de organização \textbf{baseada} em normas de conduta responsáveis por sua especificidade cultural.
Na construção de sua vida em sociedade, o indivíduo estabelece relações e interações sociais com outros indivíduos, constrói sua \textbf{percepção} de mundo, atribui significados ao mundo ao seu redor, interfere e transforma a natureza, produz conhecimento e saberes, com base em alguns procedimentos cognitivos próprios, fruto de suas tradições tanto físico material como simbólico-culturais.
A forma como diferentes sociedades estruturam e organizam o espaço físico-territorial e suas atividades econômicas permite, por exemplo, reconhecer os diversos modos como essas sociedades estabelecem suas relações com a natureza, incluindo -se os problemas ambientais resultantes dessas interferências.
Aprender a \textbf{viver} em sociedade significa, então, submeter -se a processos de socialização, ou seja, processos de incorporação e internalização de valores, papéis e identidades.
Portanto, a sociedade como teia de relações é fundamental para apreender o modo como as ações dos indivíduos \textbf{configuram} o mundo em \textbf{que} \textbf{vivem}, ao mesmo tempo em \textbf{que} constroem \textbf{uma} identidade coletiva \textbf{que} lhes permite se pensar como Nós diante do Outro ( ou Outros de referência ). c ) Sobre o conceito de Cultura : As transformações na ação das pessoas são mediadas pela cultura.
Em sua etimologia latina, a cultura \textbf{remete} à ação de cultivar saberes, práticas e costumes em \textbf{um} determinado grupo.
Na tradição metafísica, a cultura foi apresentada em oposição à natureza.
Atualmente, as Ciências Humanas compreendem a cultura a partir de contribuições de diferentes campos do saber.
O caráter polissêmico da cultura permite compreender o modo como ela se apresenta a partir de códigos de comunicação e comportamento, a partir de símbolos e artefatos e como parte da produção, circulação e consumo de sistemas de identificação cultural \textbf{que} se manifestam na vida social.
As pessoas estão inseridas em culturas ( urbana, rural, erudita, de massas, \textbf{popular} etc. ) e, dessa forma, são produtoras e produto das transformações culturais e sociais de seu tempo. d ) Sobre o conceito de Ética : O entrelaçamento entre questões sociais, culturais e individuais permite aprofundar, no Ensino Médio, a discussão sobre a ética.
Para tanto, os estudantes devem dialogar sobre noções básicas como o respeito, a convivência e o bem comum em situações concretas.
A ética pressupõe a avaliação de posturas e a tomada de posição em defesa dos direitos humanos, a identificação do bem comum e o \textbf{estímulo} ao respeito e ao acolhimento às diferenças entre pessoas e povos, tendo em vista a promoção do convívio social e o respeito universal às pessoas, ao bem público e à coletividade ( Ibidem, 2017, p. 555 ).
Por último, a quarta \textbf{categoria} apresentada, \textbf{que} está relacionada com as outras citadas anteriormente, e \textbf{que} é fundamental para o bom desenvolvimento das cidades-Estados é a \textbf{categoria} sobre Política e Trabalho.
Essa \textbf{categoria}, assim como as demais, tem a sua importância para a área de CHSA, \textbf{uma} vez \textbf{que} toda vida em sociedade \textbf{diz} respeito às práticas individuais e coletivas \textbf{que} são mediadas e cultivadas pela política e pelo trabalho dos cidadãos e dos povos \textbf{que} habitam aquela localidade.
A política, como sabemos, não é somente \textbf{uma} temática abordada pelo pensamento filosófico, mas ela vem a ser a gênese de toda a filosofia.
A política, \textbf{que} iniciou lá com os gregos antigos, para ser mais específico, com os sofistas – depois passando por Sócrates, Platão e Aristóteles, até chegar aos dias \textbf{atuais} – tinha como objetivo preparar os jovens aristocratas para virem a exercer a vida pública.
Gerir o destino da cidade não era e nem é \textbf{uma} \textbf{tarefa} fácil.
O governante deve estar preparado para exercer a vida política, pois a vida pública \textbf{diz} respeito ao bem comum, \textbf{diz} respeito à felicidade de todos.
O verdadeiro governante deve priorizar sempre o coletivo, o todo, deixando o individual em segundo plano.
Sendo assim, a política deve ser compreendida como \textbf{uma} ação social de inclusão de todos os indivíduos na sociedade.
Fazer parte da pólis - cidade-Estado - é \textbf{um} direito assegurado a todos os cidadãos.
Além da política, \textbf{um} outro fator fundamental para a sociedade é o trabalho.
Como o trabalho faz parte da vida política, todos os cidadãos, em suas respectivas funções, deveriam buscar o melhor para a sociedade, e nunca o melhor para si mesmo, ou para os seus...
O trabalho é fundamental para o desenvolvimento da economia e da sociedade.
Compreendemos \textbf{que} a vida do \textbf{homem} gira em torno do trabalho : o \textbf{homem} sempre \textbf{necessitou} do trabalho para garantir sua sobrevivência e a sobrevivência dos seus.
Nos dias \textbf{atuais}, numa visão capitalista, podemos compreender \textbf{que} o trabalho do \textbf{homem} é \textbf{um} produto \textbf{que} está à venda para a sociedade.
Sendo assim, o trabalho passa a ser \textbf{uma} mercadoria \textbf{que} está à venda.
Mas como o trabalho é desenvolvido pelo corpo humano, podemos \textbf{dizer}, de \textbf{uma} forma geral, \textbf{que} o corpo humano está à venda : vende -se aquilo \textbf{que} o corpo humano pode produzir para o bem-estar social.
O trabalhador exerce \textbf{uma} função social importante, pois o seu trabalho ajuda a desenvolver a sociedade. \textbf{Viver} em sociedade é \textbf{um} direito constitucional : todos têm o direito de “ \textbf{viver} bem ” em sociedade, e deveria ser dever do Estado garantir o bem-estar de todos.
A vida em sociedade deve ser sempre a prioridade do governante e do governado.
Desta forma, \textbf{vejamos} o \textbf{que} o documento curricular de CHSA \textbf{diz} a respeito dessa \textbf{categoria}.
A política ocupa posição de \textbf{centralidade} nas Ciências Humanas.
As discussões em torno do bem comum, dos regimes políticos e das formas de organização em sociedade, as lógicas de poder estabelecidas em diferentes grupos, a micropolítica, as \textbf{teorias} em torno do Estado e suas estratégias de legitimação e a tecnologia interferindo nas formas de organização da sociedade são alguns dos temas \textbf{que} estimulam a produção de saberes nessa área ( Ibidem, p. 556 ).
A política está na origem do pensamento filosófico.
Na Grécia Antiga, o exercício da argumentação e a discussão sobre os destinos das cidades e suas leis estimularam a retórica e a abstração como práticas necessárias para o debate em torno do bem comum.
Esse exercício permitiu ao cidadão da pólis compreender a política como produção humana capaz de favorecer as relações entre pessoas e povos e, ao mesmo tempo, desenvolver a crítica a mecanismos políticos como a demagogia e a manipulação do interesse público.
A política, em sua origem grega, foi o instrumento utilizado para combater os autoritarismos, as tiranias, os terrores, as violências e as múltiplas formas de destruição da vida pública ( Ibidem, p. 556 ).
No mundo \textbf{contemporâneo}, essas questões observadas tanto em escala local como global ganham maior visibilidade na Geopolítica, pois ela enuncia os conflitos planetários entre pessoas, grupos, países e blocos transnacionais, desafio importante de ser \textbf{conhecido} e analisado pelos estudantes ( Ibidem, p. 556 ).
A \textbf{categoria} trabalho, por sua vez, comporta diferentes dimensões – filosófica, econômica, sociológica ou histórica : como virtude ; como forma de produzir riqueza, de dominar e de transformar a natureza ; como mercadoria ; ou como forma de alienação.
Ainda, podemos falar do trabalho como \textbf{categoria} pensada por diferentes \textbf{autores} : trabalho como valor ( Karl Marx ) ; como racionalidade capitalista ( Max Weber ) ; ou como elemento de interação do indivíduo na sociedade em suas dimensões tanto corporativa como de integração social ( Émile Durkheim ).
Seja qual for o caminho ou os caminhos escolhidos para tratar do tema, é importante destacar a relação sujeito/trabalho e toda a sua rede de relações sociais ( Ibidem, 2017, p. 567-568 ). \textbf{Uma} vez \textbf{que} as \textbf{categorias} são desenvolvidas e aprofundadas no Ensino Médio, o estudante passa a compreender o quão importante é o conhecimento heterogêneo, ou seja, \textbf{um} conhecimento mais amplo e aprofundado do mundo, das relações entre seres \textbf{humanos}, das vivências cotidianas com o diferente, das diversas culturas e classes sociais, da forma como a sociedade é estruturada, das leis, da ordem, da responsabilidade por si mesmo e pelo outro, dos múltiplos papéis do \textbf{homem} na sociedade como \textbf{um} \textbf{sujeito} atuante num projeto político \textbf{que} envolve o todo.

% matched lemmas: atual, autor, basear, bastante, cartográfico, categoria, centralidade, cientista, configurar, conhecido, conseguinte, contemporaneidade, contemporâneo, dizer, enquanto, estímulo, fronteira, historicamente, hoje, homem, humanos, interdisciplinaridade, justo, necessitar, percepção, pesquisador, popular, posicionamento, que, remeter, simples, sujeito, tarefa, teoria, teórico, transdisciplinaridade, um, uma, ver, viver
\end{textsample}
